\documentclass[11pt]{article}
\usepackage[margin=1in]{geometry}
\usepackage{graphicx}
\usepackage{booktabs}
\usepackage{amsmath}
\usepackage{microtype}
\usepackage{xurl}
\setlength{\emergencystretch}{3em}
\usepackage{caption}
\usepackage{float}

\title{RE Extension}
\author{}
\date{}

\begin{document}
\maketitle

\section{RE Extension}

This section studies a rational-expectations extension of the demographic transition experiment in
\textit{NIMBYism and the Housing Supply}. The benchmark quantitative exercises in the paper treat
future house prices using the model's original forecasting rule. The natural extension is to ask
whether the main housing-market transition changes materially when households instead form
expectations over an internally consistent future price path. We keep the exercise deliberately
narrow. In particular, we abstract from coalition formation and bargaining and focus only on the
no-politics transition problem, so that the expectations channel can be isolated from the political
aggregation channel.

The extension solves a finite-horizon transition problem on the observed U.S. demographic path from
2010 to 2018. Given a guessed future price path, households solve their dynamic problem backward,
the cross-sectional distribution is simulated forward, and the implied market-clearing future price
path is recovered from housing demand and supply. A rational-expectations equilibrium requires the
guessed path to coincide with the implied path. This turns the problem into a high-dimensional
backward-forward fixed point in which the object to be solved is an entire future sequence of house
prices rather than a scalar steady state.

\subsection*{Continuation design}

Because the full rational-expectations transition is numerically demanding, we evaluate the
extension through a coarse continuation design. Let $\lambda \in [0,1]$ scale the demographic
transition around the 2010 baseline age structure. At $\lambda = 0$, every period uses the
baseline demographic profile. At $\lambda = 1$, the model uses the full 2010--2018 demographic
path. Intermediate values interpolate linearly between these two objects. Each step in the ladder
is warm-started from the previous step's final price path. The continuation ladder therefore asks
whether the full-shock case is connected smoothly to nearby easier problems or whether rational
expectations reveal a distinct transition branch.

\begin{figure}[H]
\centering
\includegraphics[width=0.98\textwidth]{figures/re_lambda_continuation.pdf}
\caption{Rational-expectations continuation across demographic shock sizes. The left panel plots the
final rational-expectations price path for $\lambda \in \{0,0.5,1\}$ as log deviations from the
baseline price level. The right panel reports the residual norm and the maximum pointwise pricing
gap along the same continuation ladder. The key pattern is smoothness: as the full demographic
transition is turned on, the rational-expectations solution moves gradually rather than jumping to a
qualitatively different branch.}
\end{figure}

Figure 1 delivers the main message of the continuation exercise. The full-shock
rational-expectations path is not dramatically different from the paths obtained under smaller
demographic shifts. The continuation statistics also improve monotonically in residual norm as
$\lambda$ rises, while the largest pointwise pricing error returns to the familiar difficult region
near the full-shock case. At low values of $\lambda$, the hardest period sits late in the path. As
the full demographic transition is restored, the bottleneck shifts back toward the third model
period. This is useful evidence for interpretation. It suggests that the extension is not hiding a
large, economically distinct rational-expectations equilibrium; rather, the same localized
bottleneck becomes harder to clear as the full demographic transition is switched on.

\subsection*{Benchmark and challenger RE solutions}

Within the rational-expectations extension itself, we compare two full-shock solutions. The
benchmark solution uses the most stable global-control update rule. The challenger relaxes the
candidate-selection rule so that the solver is more willing to trade average residual fit against
local worst-gap performance. This comparison is informative because it shows whether the economic
transition path is sensitive to the numerical details of the update rule.

\begin{figure}[H]
\centering
\includegraphics[width=0.98\textwidth]{figures/re_solver_comparison.pdf}
\caption{Benchmark and challenger rational-expectations solutions under the full demographic
transition. The left panel plots the final price paths implied by the benchmark control rule and the
best relaxed challenger. The right panel plots the corresponding raw log-price residuals. The
challenger slightly improves average residual fit, but the two price paths remain close and the
largest residual remains concentrated in a narrow set of periods.}
\end{figure}

Figure 2 shows that the economic path is much more stable than the raw numerical diagnostics might
suggest. The relaxed challenger slightly lowers the overall residual norm and keeps nontrivial
updates alive for one additional iteration, but it does so by worsening the largest local pricing
gap. The two final price paths are nevertheless very similar. The differences are concentrated in a
small set of later periods rather than spread across the whole transition. This reinforces the
interpretation from the continuation exercise: rational expectations appear to matter at the margin,
but the extension does not reveal a qualitatively different transition in this calibration.

\subsection*{Heuristic comparison to the paper benchmark}

The question for the paper is not only whether the rational-expectations solver converges cleanly,
but also whether allowing households to internalize the future price path seems to bring the NIMBY
mechanism forward in time relative to the published benchmark. The clean apples-to-apples comparison
is not available from a shared export file, because the published benchmark retains the political
aggregation channel while the present extension shuts politics down to isolate expectations.
Nevertheless, the saved benchmark figure used in the original project is informative enough to
construct a directional comparison over the overlapping 2010--2018 window. Figure 3 therefore
normalizes both series at 2010 and overlays the full-shock rational-expectations path on a
digitized version of the local benchmark figure.

\begin{figure}[H]
\centering
\includegraphics[width=0.82\textwidth]{figures/re_vs_paper_benchmark.pdf}
\caption{Heuristic comparison between the published benchmark and the rational-expectations
extension over 2010--2018. The black line is digitized from the local saved benchmark figure
\texttt{Baseline.jpg} in the original project workspace and normalized at 2010. The orange dashed
line is the full-shock rational-expectations no-politics transition, also normalized at 2010. The
comparison is intentionally directional rather than exact because the two environments differ in the
presence of coalition politics. The main pattern is that the benchmark rises steadily through the
window, while the rational-expectations extension remains nearly flat.}
\end{figure}

Figure 3 sharpens the interpretation. The current rational-expectations extension does not by
itself appear to ``bring NIMBY forward.'' Over the shared 2010--2018 window, the benchmark path
rises materially while the no-politics rational-expectations path barely moves relative to its 2010
level. The most reasonable reading is therefore not that rational expectations generate an earlier
version of the benchmark transition. Instead, the benchmark's stronger short-run appreciation seems
to be tied to the broader political benchmark environment rather than to model-consistent price
expectations alone.

\subsection*{Why the full RE transition is computationally hard}

The computational difficulty in this extension is structural rather than a matter of insufficient
hardware. First, the equilibrium object is an entire future price path. Second, each candidate
price path requires a full backward solution of household decisions and a forward simulation of the
distribution. Third, the pricing errors are not uniform across periods: they are sharply localized.
The continuation exercise shows that the hard region returns to the third model period once the full
demographic transition is restored. Finally, alternative update rules reveal a real ranking
problem. A candidate can improve the overall residual norm while worsening the largest pointwise
pricing error. That is why a brute-force increase in computation is not enough by itself. The main
limitation is the algorithmic difficulty of moving the entire path while respecting the localized
equilibrium condition in the difficult periods.

\subsection*{Interpretation}

Taken together, the extension points to a cautious but economically useful conclusion. Allowing for
rational expectations along the demographic transition does not overturn the benchmark message of
the paper. Instead, it generates modest and localized adjustments to the housing price path. The
full-shock rational-expectations case reconnects smoothly to nearby easier continuation problems,
and the benchmark and challenger rational-expectations paths remain close even when the numerical
ranking criteria differ. For the purposes of the paper, this is enough to justify a qualitative
reading of the extension: rational expectations do not appear to create a fundamentally different
housing-market transition under the present calibration, although they do sharpen the numerical
fixed-point problem in a small set of periods.

\end{document}
